\section{Synchronized Planarity with complete-sector constraints}
\label{cse:section}

The central change from the two-anchor engine is to retain an entire sector
category. A vertex may have many marked directions in an order prescribed
up to reflection, while its other directions remain freely insertable in
every sector. The only relations involving large categories will be
injective correspondences. Binary quantities may satisfy arbitrary affine
equations. This restriction is what separates the present language from
the NAE-hard language of arbitrary multi-gap observations.

\subsection{The instance and its preservation contract}

Let $G$ be a loopless multigraph with the ordinary P-vertices, Q-vertices,
Q-cells, and matching of pipes of Synchronized Planarity. A pipe
$(u,v,\phi)$ requires
\[
 \rho_v=\rev\bigl(\phi(\rho_u)\bigr).
\]
A Q-vertex has a prescribed reference rotation and uses it or its reverse
according to its cell's Boolean variable. At degrees zero, one, and two,
the rotational requirement is vacuous, while the cell variable is retained.

\begin{definition}[Complete-sector instance]
\label{cse:definition}
Mark a subgraph $L\subseteq G$ satisfying the following conditions.
\begin{enumerate}
\item Every marked edge belongs to a simple marked cycle.
\item At each marked P-vertex, all marked incident edges belong to one
      biconnected block of $L$.
\item A pipe is either off at both ends or maps the complete marked dart
      set at one end onto the complete marked dart set at the other.
\end{enumerate}
At each marked P-vertex $v$, prescribe a reference cyclic list
\[
 A_v=(a_0,\ldots,a_{k_v-1})
\]
of all its marked darts and a Boolean frame variable $f_v$. Require
\[
 \rho_v|_{A_v}=A_v\quad\text{if }f_v=0,
 \qquad
 \rho_v|_{A_v}=\rev(A_v)\quad\text{if }f_v=1.
\]
Here $k_v\ge2$. For $k_v=2$ the order condition is vacuous; $f_v$ remains
an explicit gauge for the two sector names.

For an unmarked dart $d$ at $v$, define its category $c_v(d)$ to be the
unique $j\in\{0,\ldots,k_v-1\}$ such that
\[
 \begin{cases}
 d\in(a_j,a_{j+1})_{\rho_v},&f_v=0,\\
 d\in(a_{j+1},a_j)_{\rho_v},&f_v=1,
 \end{cases}
\]
where indices are cyclic. The two sector names remain distinct when
$k_v=2$.

Variables have explicitly listed finite domains. An observer has the form
$x=c_v(d)$, with the corresponding $k_v$-element domain. A relation between
two finite-domain variables is allowed if it is a \emph{partial bijection}:
it contains at most one pair with any prescribed first coordinate and at
most one pair with any prescribed second coordinate. Unary domain
restrictions are allowed. Boolean variables are designated two-element
variables with values $0,1$. Arbitrary affine equations over $\Ftwo$ may
involve these Boolean variables, including Q-cell variables, frame
variables, and two-sector variables. In particular, every observer
$c_v(d)$ at a vertex with $k_v=2$, every two-valued category variable
introduced by the operations below (such as a detached-component variable
$h_T$ in the Join when $k=2$, and the two-direction source variables
$\xi_{\mu,e}$ of the pipe-free stage), and every selector $b$ is designated
Boolean; its values $0,1$ are its two sector names. A noninjective Boolean
predicate of a larger category is not permitted as an affine variable.
\end{definition}

A witness comprises a planar rotation system and values of all variables
that satisfy the ordinary pipes and Q-cells, all frame and observer
conditions, all partial bijections, and all affine equations. An operation
preserves the instance if it preserves the exact set of assignments of
\emph{every} old variable and has constructive witness maps in both
directions. This includes gauge variables introduced by earlier
operations. No operation may silently replace their values.

An active graph observer is interned once for each pair $(v,d)$.
Repeated input observer identifiers are joined to it by identity partial
bijections. Relocating an old variable deletes its old active graph
observation, retains the variable, and adds one indicated transport
relation to the newly interned representative. Thus repeated earlier
identifiers do not generate repeated graph observations in later steps.

\begin{theorem}[Complete-sector engine]
\label{cse:main}
Let $M$ be one plus the explicit input size, including all domain lists,
relation tables, reference lists, and affine rows. Feasibility of a
complete-sector instance is decidable in deterministic
$O((M+1)^{20})$ time. Within the same bound, the algorithm constructs a witness.  The
large-domain variables are not assumed to admit an affine encoding.
\end{theorem}

Conditions 1--3 of Definition~\ref{cse:definition} are hypotheses on the
input: the engine does not test them, and Theorem~\ref{cse:main} applies
to instances that satisfy them. For compiled instances they are
established in Section~\ref{gen:marking}.

The exponent is deliberately conservative. Its purpose is to give an
explicit bound while retaining simple implementations of every table and
source-tree operation.

\subsection{The finite-domain algebra}

\begin{lemma}[Partial bijections and Boolean equations]
\label{cse:algebra}
A system of partial-bijection relations, unary domain restrictions, and
affine equations on designated Boolean variables is polynomial-time
solvable. Its feasible relation on those Boolean variables is affine, and
every feasible Boolean assignment can be extended constructively.
\end{lemma}
\begin{proof}
Consider one connected component of the graph whose edges are the
partial-bijection relations. Fix a root variable. For each possible root
value, propagate its unique image or inverse image across every edge.
Discard the root value on encountering an undefined image, an excluded
domain value, or inconsistent values at a vertex. A surviving root value
determines one complete assignment of the component. Conversely, every
solution is obtained this way.

If the component contains a Boolean variable, choose one as root. There
are at most two surviving assignments. Zero assignments reject the
instance. One assignment fixes every Boolean variable in the component.
If two survive, each other Boolean variable has the form $t\oplus b$,
where $t$ is the root value and $b$ is a constant. Indeed, a path composes
injections, so its two different root values cannot produce the same
value at the other end.

A component with no Boolean variable occurs in no affine equation. It
can independently choose any surviving root value. Substitute the affine
description of the other components into the given affine equations and
solve the resulting linear system over $\Ftwo$. Back substitution chooses
the finite-domain values. Gaussian elimination also gives the claimed
projection. All enumerated root domains and relation tables are explicit
parts of the input, so the procedure takes polynomial time.
\end{proof}

The injectivity assumption is indispensable to this argument. If a large
category could be mapped to a Boolean interval-membership predicate,
different category values could agree at one Boolean variable while
disagreeing at another; the preceding two-assignment conclusion would no
longer follow.

\subsection{Pipe transfer and source structure}

Suppose a pipe maps $k\ge3$ marked darts. Its map must send the first
reference cyclic list either to the second reference list or to its
reverse, allowing a cyclic shift. Otherwise no admissible pair of marked
orders satisfies the pipe, and the instance is rejected. If the reference
orders are preserved, the pipe requires $f_v=f_u\oplus1$; if they are
reversed, it requires $f_v=f_u$. Under this relation it induces a fixed
bijection $\sigma_{uv}$ of reference sector names, and
\begin{equation}
\label{cse:large-transfer}
 c_v(\phi(d))=\sigma_{uv}\bigl(c_u(d)\bigr).
\end{equation}
These statements follow by applying the reference map and then reversing
the full cyclic order. Retain the forced frame equation if the pipe is
subsequently removed.

For two marked darts the reference cyclic order does not determine frame
parity. Choose their stored ordered representatives. The ordinary
two-gap reads satisfy $g_u\oplus g_v=e$ for a reference-dependent
constant. Since a framed category is its underlying gap read, its frame
bit, and a fixed labeling offset added modulo two, the exact transfer is
\begin{equation}
\label{cse:two-transfer}
 c_u(d)\oplus c_v(\phi(d))\oplus f_u\oplus f_v=e_{uv}.
\end{equation}
No equation between $f_u$ and $f_v$ alone is inferred in this case.

Replace each Q-vertex of degree at least three by its ordinary reference
wheel. Lift marked edges through their edge segments and spokes; no rim
edge is marked. Let $G^W,L^W$ denote the resulting graph and marking. The
ordinary wheel witness maps preserve rotations of old P-vertices. Every
port has marked degree at most two, while a hub may meet several marked
blocks.

\begin{lemma}[One complete source anchor set]
\label{cse:source}
No wheel hub is a P-source pole. A P-source with at least four directions
(skeleton edges) has original P-vertices at both poles. If a P-source displays marked cycles,
all those cycles belong to one block of $L^W$, and the set of source
directions used by that block is the same at both poles. A source with at
least three marked directions has original P-vertices at both poles.
\end{lemma}
\begin{proof}
If a hub $q$ and another vertex $w$ were P-source poles, the rim minus
$w$ would connect every port other than possibly $w$ while avoiding both
poles. All those spokes lie in one expansion, and only a possible spoke
$qw$ lies in another. There are at most two directions, contradicting the
fact that a P-source has at least three. A wheel port cannot be a pole of a
P-source with at least four directions: its other pole is not the hub, and after
deleting both poles its spoke and at least one rim edge remain joined
through the hub. These two incident edges enter one expansion, leaving
at most three directions from its at most four graph edges. Unwheeled
Q-vertices have degree at most two. This proves the claim about four
directions.

Every displayed marked cycle contains both source poles. Two distinct
biconnected blocks cannot share both poles, so the displayed cycles
belong to one marked block $B$. At either pole all marked incident edges
belong to this block: at an original P-vertex by the second invariant, and
at a wheel port because a port has at most two marked edges and a marked
cycle through one of them, which exists by the first invariant, uses the
other. If one such edge
enters an expansion at a pole $p$, connectivity of $B-p$ connects its
other end to the other pole $q$ through that expansion: the only exits
from an expansion are its poles. Hence the expansion has a marked edge
at $q$ as well. The symmetric argument identifies the same complete
direction set at both poles. Finally a wheel port has at most two marked
edges, and a hub is excluded. Thus three or more marked directions
require original P-poles.
\end{proof}

We use the ordinary SPQR description of planar embeddings: independent
R-source reflections and P-source permutations, with opposite orders at
the two poles. The source-labelled embedding tree at a vertex has one
leaf per incident dart. Its frontier gives the vertex rotation, and the
orientation of three leaves is the orientation of their directions at
their unique median. These are the ordinary embedding-tree facts used by
the BFR operations~\cite{bfr}.

\subsection{Small vertices and encapsulation}
\label{cse:small}

\paragraph{ConvertSmall.}
Convert a P-vertex of degree at most three, and its pipe mate when
present, exactly as in the ordinary algorithm. As in BFR~\cite[Section
3.1.1]{bfr}, the conversion is applied whenever such a vertex appears: at
the start, and again after every operation of the schedule in
Section~\ref{cse:schedule}, once that operation's relocations have been
made (an encapsulate-and-join step counts as one operation). Thus every
P-vertex has degree at least four whenever an operation is applied. Use the new Q-cell
variable $g$ and the ordinary reversed matched reference rotations.
With three marked darts there is no private observation, and the old
frame condition is the affine equation $f_v\oplus g=b$. With two marked
darts, an observed third dart satisfies
\[
 c_v(d)=g\oplus f_v\oplus b,
\]
where $b$ compares the chosen references. Off vertices have neither
condition. Keeping the old variables and adding these equations preserves
all assignments in both directions. Marks are unchanged; their block
condition no longer needs to be checked at a converted Q-vertex.

\paragraph{Encapsulation.}
Let $u$ be a P-cutvertex. For every component $Y$ of $G-u$, create the
ordinary residual cap and star ray, linked by the cap pipe. Mark the two
replacement edges exactly when the old edge at $u$ was marked.

If $u$ is marked, all its marked darts enter one component of $G-u$,
because their marked block remains connected after deleting $u$. Thus
one cap and one ray have the same complete marked set as $u$; all other
new caps and rays are off. The marked part of the star is a parallel
bond with $k_u$ edges. The residual marked block is the old block with
$u$ renamed as its cap, and is biconnected. Every other marked block is
unchanged. The new pipe maps complete marked sets.

Give the marked cap the old reference list, and the marked ray its
reverse. Introduce fresh frame variables at both and equate both with
$f_u$, including when $k_u=2$. This does not restrict any old variable:
the fresh gauges can always be chosen in this way. The ordinary split
map preserves all old rotations, and its inverse uses the reversed cap
correspondence. Thus every old category and frame is preserved. The
reference choices make the new frame and category conditions exact.

BFR allow the two ends $u,v$ of the pipe to lie in one component of $G$,
possibly adjacent or on one marked block. Encapsulate $u$ first. Then $v$
is still a cutvertex of its new piece: a component of $G-v$ avoiding $u$
has no edge to $u$, hence none to the cap of $u$, and remains a component
after deleting $v$ from the piece. A marked block through $u$ and $v$
becomes the residual block of $u$, with $u$ renamed as its cap; after
encapsulating $v$ it is a residual block with two caps. The arguments
above apply to each end separately.

\subsection{Joining stars with many marked edges}

Join the two encapsulated stars with centers $u,v$ along their center
pipe. The ray graph $K$ is bipartite; its edge $\ell_d$ corresponds to a
matched pair of old center darts. Every ray rotation retains its identity.
If the centers are marked, the $k$ marked edges form a parallel bond
between the two marked rays $p,q$.

Keep the ray reference lists and frame variables supplied by
encapsulation. In particular,
\[
 A_p=\rev(A_u),\qquad f_p=f_u,
 \qquad A_q=\rev(A_v),\qquad f_q=f_v,
\]
under the appropriate dart identifications. Do not overwrite $f_q$ with
$f_u$: for two anchors these may be different old values. Measure global
ray-graph sectors in $p$'s frame. The star bonds induce fixed reference
sector bijections from $u$ to $p$ and from $v$ to $q$. Transfer a local
$q$-category to $p$ by a fixed bijection when $k\ge3$, using the retained
old pipe-frame equation; for $k=2$ use
\[
 c_q\oplus f_q\oplus f_p\oplus b
\]
with the computed reference constant $b$.

\begin{lemma}[Complete-sector Join]
\label{cse:join}
Every old center observation can be relocated to a category at $p$ or
$q$, or to one fresh category for a detached component of $K$, with
partial-bijection relations and affine Boolean equations. The resulting
instance preserves every old assignment and has constructive witnesses
in both directions.
\end{lemma}
\begin{proof}
In a forward witness, all edges of an ordinary ray have the same global
sector label. Otherwise that ray's star bond would interleave with two
edges of the marked star bond. At $p$ and $q$, opposite bond orders give
exactly the transfer formulas above. Labels are therefore constant along
every connected piece avoiding $p,q$.

Let $K_0$ contain the marked bond. If an observed edge is in $K_0$, follow
a fixed path from it to $\{p,q\}$, first meeting that set at the last
vertex. Its label is the category of the last edge there, converted to
$p$'s frame. This uses a fixed sector bijection for $k\ge3$ and an affine
formula for $k=2$. For every other component $T$ introduce one arbitrary
category $h_T$ in the $k$-element $p$-sector domain. Relocate the old
$u$- and $v$-observations through their retained star-bond transfers.

Conversely, the ordinary BFR separating-curve inverse splits each planar
embedded component of $K$ into two stars with opposite matched center
orders and unchanged ray rotations~\cite{bfr}. Keep the center order
obtained from $K_0$. For a detached component $T$, convert $h_T$ through
the fixed $p$-to-$u$ sector bijection and insert a linear representative
of its entire center order in that $u$-sector. If its name is $j$, insert
after $a_j$ when $f_u=0$, or after $a_{j+1}$ when $f_u=1$. Use the reversed
matched complete order at $v$.

Each component is a contiguous excursion, preserving the noncrossing
ray-bond partitions of both stars. Its sector is exactly the selected
$h_T$; later insertions in the same sector do not change this. The
restriction to $K_0$ is unchanged, and path-label propagation identifies
every recovered category there with its relocated value. Transfer recovers
the observations at the other center while retaining both old gauges.
All ray rotations, cap pipes, and residual-piece constraints survive.
\end{proof}

For an off center pipe, all star edges are unmarked, and there are no
center frames or observations. Use the ordinary join and concatenate its
component orders in the inverse. Marks elsewhere are not removed: every
marked cycle lies in a residual piece and is unchanged. The cap pipes are
off at their newly created ends. Thus the off case also preserves the
whole marked instance.

\subsection{Propagation and injective contour maps}

Let $u$ be a matched block-vertex with a nontrivial ordinary embedding
tree $T$, computed in $G^W$. Let $S$ be the minimal subtree spanning its
marked leaves. A cyclic reference order of these leaves is a frontier
of the unlabelled tree $S$ precisely when every tree-edge shore has its
leaf set in one cyclic interval. To see sufficiency, the reference order
determines the cyclic order of the branches at every branching vertex;
tracing the contour with those branch orders gives the prescribed leaf
order. Necessity follows because a branch occupies one interval in every
plane-tree frontier. At an inner Q-node the induced order of three or
more marked branches must additionally match its prescribed order or its
reverse. Check these conditions explicitly and reject if any fails.

\begin{lemma}[Injective sector transport in a tree]
\label{cse:contour}
Fix a reference plane embedding of $S$ with the prescribed marked-leaf
order. At an inner vertex $m$ with $h$ marked branches, its $h$ reference
corners map injectively to the intervals between consecutive marked
leaves. If a private leaf attaches to $S$ at $m$, its full sector at the
old vertex is the image of its local sector at $m$. The map remains
unchanged when the reference tree and local frames are reflected together.
\end{lemma}
\begin{proof}
Follow the tree contour through a corner at $m$. It passes from the last
leaf of one incident branch to the first leaf of the next. These are
consecutive leaves in the full frontier. Two different corners at $m$
lie between different consecutive branch intervals, and hence determine
different leaf intervals. A private branch hanging at $m$ is traversed
exactly at its attachment corner, proving the first identity. Reversing
the whole reference reverses both traversals while leaving each named
reference corner and its named reference leaf interval corresponding.
The same argument applies to two branches, whose two corners are distinct.
\end{proof}

Perform the ordinary PropagatePQ replacement with two copies of $T$,
retaining leaves as degree-two subdivisions. Each such subdivision vertex
is an unmatched Q-vertex of degree two in a fresh singleton Q-cell: its
rotation requirement is vacuous, and it carries no frame and no observer.
This departs from the BFR graph operation only by these subdivisions; the
size estimate of Section~\ref{cse:schedule} accounts for them. Mark the
spanning subtree
in each copy. Crucially, carry out the reference construction separately
using each endpoint's own old reference and frame. New P-nodes in the
$u$-copy use their induced reference marked-branch lists and fresh frames
equal to $f_u$; the corresponding nodes in the $v$-copy use its own
induced references and frames equal to $f_v$. These two gauges are not
silently identified when there are two marked leaves.

Corresponding inner P-nodes receive the ordinary pipes. Corresponding
inner Q-nodes receive a shared fresh cell and reversed full reference
orders. At a Q-node with at least three marked branches, enforce the
affine equation between its native Q-flip and the old frame that makes
the induced marked order correct. With two marked branches impose no
frame equation.

For a private leaf $d$, let $m$ be its first point on $S$, and let
$\delta$ be the direction toward $d$. If $m$ is P, relocate an old
category $x$ by
\[
 x=\iota_m\bigl(c_m(\delta)\bigr),
\]
where $\iota_m$ is the injection of Lemma~\ref{cse:contour}; this is a
partial-bijection relation. If $m$ is Q with three or more marked
branches, its category is constant after its frame equation, and $x$ is
pinned to the corresponding image. If $m$ is Q with two marked branches,
introduce a Boolean variable
\[
 b=g_m\oplus f_u\oplus b_0
\]
in the $u$-copy, using $f_v$ in the $v$-copy, and impose
$x=\iota_m(b)$. The constant $b_0$ compares the full Q-reference with
the two chosen local corners.

The ordinary forward map expands center rotations into tree rotations;
the inverse contracts the trees. The frontier test enforces precisely the
old marked orders, and Lemma~\ref{cse:contour} enforces every old category
in both directions. Corresponding marked subtrees map to each other under
the old pipe, so new pipes align on whole marked sets. Their transfer
uses \eqref{cse:large-transfer} or \eqref{cse:two-transfer} as appropriate.

The marked-block invariant is retained. Replacing a vertex of a
biconnected marked block by the minimal tree spanning all its incident
marked leaves leaves a biconnected graph. After deleting a new tree
vertex, every remaining branch contains a marked leaf, and the old
marked block minus the old vertex connects all those leaves outside the
tree. Deleting an old vertex preserves connectivity by contracting the
new tree and using biconnectivity of the old block. This also shows that
all inserted marked edges remain cyclic. Off pipes insert no marked
edge and have no frame or observer updates.

\subsection{Superblocks and SimplifyMatching}

Suppose $u$ is a matched P-block-vertex with a trivial embedding tree and
degree $D\ge4$. Let $w$ be its partner pole. Its source expansions have
one dart $e_j$ at $u$ and a nonempty interval $I_j$ in the relevant block
rotation at $w$. Define $S_j$ to start with $I_j$ and include all
other-block insertions up to the next such interval. Changing the order
at $u$ and concatenating these unchanged linear superblocks in the
opposite expansion order at $w$ preserves planarity and all other
rotations. We prove this first.

\begin{lemma}[Superblock rearrangement]
\label{cse:superblock-planar}
Let $\rho$ be a planar rotation system of $G^W$, and let $\tau$ be a cyclic
order of the darts at $u$. Define $\rho^\tau$ by $\rho^\tau_u=\tau$,
$\rho^\tau_x=\rho_x$ for $x\ne u,w$, and let $\rho^\tau_w$ be the cyclic
concatenation of the superblocks, each with its linear order, in which
$S_k$ directly follows $S_j$ exactly when $e_j$ directly follows $e_k$ in
$\tau$. Then $\rho^\tau$ is planar, and $\rho^{\rho_u}=\rho$.
\end{lemma}
\begin{proof}
Let $B$ be the block containing $u$; it contains all darts at $u$, since
$u$ is not a cutvertex. Its embedding tree at $u$ is the single P-node
$\mu$ with poles $u,w$, and by Lemma~\ref{cse:source} both poles are
original P-vertices. In the embedding $\mathcal E$ of $B$ induced by
$\rho$, the order of the $e_j$ at $u$ is the reverse of the order of the
intervals $I_j$ at $w$, so $\rho^{\rho_u}=\rho$.

Let $\mathcal E^\tau$ be the embedding of $B$ with the same SPQR choices
as $\mathcal E$, except that the skeleton order of $\mu$ at $u$ is $\tau$.
A P-skeleton order affects rotations only at its poles, so
$\mathcal E^\tau$ and $\mathcal E$ agree at every vertex other than $u,w$,
and $\mathcal E^\tau$ has rotation $\tau$ at $u$. In the embedding tree of
$B$ at $w$ only the rotation of the node $\mu$ changes, to the reverse of
$\tau$. The frontier concatenates the excursions at $\mu$, which are the
intervals $I_j$ in their unchanged linear orders. Hence the rotation of
$\mathcal E^\tau$ at $w$ is the restriction of $\rho^\tau_w$ to $B$.

Now consider another block $B'$ at $w$. Since $B$ has at least four darts
at $w$ and does not interleave with $B'$ in $\rho_w$, all darts of $B'$
lie in one \emph{slot} between cyclically consecutive darts $g,g'$ of $B$.
Whether $g,g'$ lie in one interval $I_j$ or $g$ ends $I_j$, the whole slot
lies in the superblock containing $g$, directly after $g$. Superblocks keep
their linear orders, so in $\rho^\tau_w$ every slot is still a contiguous
run with its old order, directly after the same dart of $B$. Consequently
$\rho^\tau_w$ restricts to the old rotation of every block other than $B$,
and no two blocks at $w$ interleave: two different slots are disjoint
intervals, and a slot contains no dart of $B$.

Apply the block-assembly fact of Section~\ref{rep:elementary} to the
component of $G^W$ containing $u$, with $\mathcal E^\tau$ on $B$, the old
rotations on all other blocks, the old merges at every cutvertex other
than $w$, and $\rho^\tau_w$ at $w$. It gives a planar rotation system,
which is $\rho^\tau$. Other components are unchanged. No wheel vertex is
$u$ or $w$, so every wheel keeps its rotation.
\end{proof}

\begin{lemma}[Complete-category superblock preservation]
\label{cse:superblocks}
If $u$ is off, an arbitrary superblock permutation preserves every marked
order and category at $w$. If $u$ is marked, a superblock permutation
preserving its marked order and every unmarked-dart category preserves
every marked order and category at $u,w$, with all old frame values.
\end{lemma}
\begin{proof}
Planarity of the rearranged rotation system is
Lemma~\ref{cse:superblock-planar}.
If $w$ is marked and $u$ is off, the marked block through $w$ avoids $u$.
If that marked block belongs to the active graph block, connectivity
after deleting $w$ puts all its marked darts in one source expansion,
hence one superblock. If it belongs to another graph block, its connected
interior lies in one face of the biconnected active block, and hence in
one corner at $w$; again its darts lie in one superblock. Its internal
order remains fixed; every other superblock is in the same complementary
marked sector. Thus all categories and the full marked order survive.

If $u$ is marked, its marked block contains $w$. By
Lemma~\ref{cse:source}, the marked superblocks at $w$ correspond exactly
to the marked darts at $u$. Their cyclic order is unchanged. Inside a
marked superblock its entire linear order, including all inserted other
blocks, remains unchanged. An unmarked superblock remains between the
same consecutive marked superblocks because its complete category at
$u$ is unchanged. This preserves every category at $w$. Frame values
are retained; since the marked orders and categories remain as before,
their defining conditions still hold.
\end{proof}

Use the ordinary BFR SimplifyMatching cases. For a pipe $(u,r,\phi)$,
delete it if $w$ is unmatched; use the self-pipe rule below if $r=w$;
or, when $w$ is a matched trivial pole with a different mate $w'$, replace
the two pipes by the ordinary shortcut
\[
 (r,w',\phi'\beta\phi^{-1}),
\]
where $\beta$ identifies opposite-pole darts. In the latter two cases
$\deg w=D$ and each $I_j$ is a singleton.

In a marked non-self case, before removal add exact transfer for
\emph{every} unmarked dart at $u$: the partial bijection
\eqref{cse:large-transfer} with its forced frame equation when $k\ge3$,
or the affine equation \eqref{cse:two-transfer} when $k=2$. In a backward
witness take
\[
 \tau=\rev\bigl(\phi^{-1}(\rho_r)\bigr)
\]
as the new order at $u$. The retained equations say exactly that $\tau$
has the old marked order, the old complete categories, and the old frame
value. Lemma~\ref{cse:superblocks} performs the required planar
rearrangement at $w$. The ordinary reversed-order composition recovers
both deleted pipes in the shortcut case. Complete marked sets correspond
under the constituent maps, so the shortcut stays aligned. In the
two-anchor case no additional gauge relation is imposed across the
implicit map $\beta$.

In an off case use the ordinary update. If $w$ is marked in the unmatched
case, Lemma~\ref{cse:superblocks} preserves its data. In the two equal-degree
cases $w$ must also be off: two marked darts in a block avoiding $u$
would have to lie in the same one-dart interval $I_j$, which is impossible.

\paragraph{Self-pipes.}
Let $\pi=\phi\beta^{-1}$ be the induced permutation on the darts at $w$.
The required order is $\pi$-invariant. A permutation admits an invariant
cyclic order precisely when all its orbits have equal length: in an
invariant order it commutes with the successor cycle and is a power of
that cycle; conversely, arranging orbit representatives and then their
successive images gives an invariant order. The off case uses this test.

In a marked case, $\pi$ preserves the anchor set. For $k\ge3$ its action
on the prescribed reference anchor cycle must be a rotation by some
$r\pmod{k}$. Reject otherwise. Its orbit length there is
$\ell=k/\gcd(k,r)$, with $\ell=1$ for $r=0$, and every orbit on all
darts must have this length. Add the bijection constraints
\begin{equation}
\label{cse:self-sector}
 c_w(\pi(d))=c_w(d)+r\pmod{k}
\end{equation}
for every unmarked dart. Necessity follows because an invariant cyclic
order carries each reference gap to that shifted gap.

For sufficiency choose an arbitrary private-dart order in one
representative of each sector orbit under $j\mapsto j+r$, and propagate
those orders by $\pi$. Each private orbit has exactly one member in each
sector of its sector orbit; this follows from \eqref{cse:self-sector} and
its length $\ell$. Use the prescribed reference anchor order for $f_w=0$
or its reverse for $f_w=1$. The completed order is invariant and preserves
every $w$-category and its frame. Recover the order at $u$ through the
bond correspondence $\beta$: here $\deg w=D$, every $I_j$ is the singleton
$\{\beta(e_j)\}$, and there are no other blocks at $w$, so the
superblocks are single darts and Lemma~\ref{cse:superblock-planar}, with
$\tau$ the reverse of $\beta^{-1}$ applied to the new order at $w$, gives
planarity. The unchanged marked order and opposite-pole correspondence
preserve all old $u$-categories as well.

For $k=2$, equal orbit lengths allow only the identity or a fixed-point-free
involution exchanging the anchors. In the latter case add
\[
 c_w(\pi(d))\oplus c_w(d)=1.
\]
Choose an arbitrary order in one sector and put its image under $\pi$ in
the other. The resulting half-turn-invariant order preserves the selected
categories. The independent gauges at the two poles require an explicit
check: for a fixed reference constant $t$, the old planar pair satisfies
\[
 c_u(d)\oplus c_w(\beta(d))\oplus f_u\oplus f_w=t.
\]
The reconstructed reversed pair satisfies the identical equation. Every
$c_w$ value and both $f_u,f_w$ are retained, so subtracting the equations
preserves every $c_u$ value. No equality between the two gauges has been
assumed.

All Simplify cases preserve every old assignment. They change neither
$G$ nor $L$ and introduce at most one complete-sector relation per dart,
plus affine frame equations.

\subsection{Solving the pipe-free instance}

When no pipe remains, test planarity of $G^W$ and compute its blocks and
SPQR sources. Retain all old variable relations. Introduce a reflection
bit $r_\mu$ for every used R-source. At each marked P-vertex $v$, all its
marked darts lie in one graph block. Form the subtree $S_v$ spanning
these leaves in its embedding tree, and perform the reference-frontier
tests used for propagation.

At a source with at least three marked directions, the projected marked
order is fixed by the frame $f_v$. For an R-source, compare it with the
native reference and add $r_\mu\oplus f_v=b$. For a P-source, choose a
canonical reference order on its complete marked direction set $A_\mu$
and a source frame $g_\mu$. All occurrences give affine equations between
$g_\mu$ and their $f_v$. Reject if any projected reference is neither the
canonical one nor its reverse. Lemma~\ref{cse:source} ensures that these
are references for the same full direction set at the two poles. These
tests and equations are necessary and sufficient for all marked-order
requirements, by the independent source choices and the tree-frontier
criterion.

Now expand each category observer $x=c_v(d)$ as follows.
\begin{enumerate}
\item If $d$ belongs to another graph block $B'$ at $v$, introduce one
      arbitrary $k_v$-valued variable $\eta_{v,B'}$ for that whole block
      and impose $x=\eta_{v,B'}$.
\item Otherwise let $m$ be the attachment of $d$'s branch to $S_v$, and
      use the injective contour map $\iota_m$.
      At an R-source with at least three marked directions, its category
      is constant because $r_\mu$ has already been synchronized with
      $f_v$.
\item At an R-source with two marked directions, use a Boolean variable
      $b=r_\mu\oplus f_v\oplus b_0$ and the partial bijection
      $x=\iota_m(b)$.
\item At a P-source with at least three marked directions, introduce one
      independent source category $z_{\mu,e}$ for every observed
      unmarked source direction $e$. Its domain is the complete reference
      sector set of $A_\mu$. Then $x$ is an injective relabeling of
      $z_{\mu,e}$, including the fixed pole/reference correspondence.
      The affine source-frame equations make the relabeling independent
      of the frame values.
\item At a P-source with two marked directions, introduce its underlying
      binary gap variable $\xi_{\mu,e}$. Use a Boolean
      $b=\xi_{\mu,e}\oplus f_v\oplus b_0$ and impose
      $x=\iota_m(b)$. The reference constant includes reversal at the
      other pole.
\end{enumerate}
Finally add the ordinary Q-wheel equations linking each nonvacuous
Q-cell flip to the reflection of the R-source containing its wheel.
Only injections involve larger domains; all remaining equations involve
Boolean variables. Lemma~\ref{cse:algebra} solves this system.

Necessity is given by the source and contour identities. For sufficiency,
choose each R-reflection from its bit. At a P-source with $h\ge3$ marked
directions, choose their order using $g_\mu$ and put each unmarked
direction in its selected reference sector, with arbitrary order within
each sector. At a source with two marked directions, use its selected
binary gaps. Unused directions may be assigned arbitrarily. The source
and contour facts guarantee every marked order and category
simultaneously. At a marked cutvertex insert every other graph block as
one contiguous excursion in its selected $\eta$ sector, preserving the
block's internal rotation. The block partitions are noncrossing, and
these insertions affect no internal source choice. Use arbitrary
noncrossing assemblies elsewhere and delete the Q-wheels. This constructs
the required planar witness.

\subsection{Schedule, size, and reconstruction}
\label{cse:schedule}

Canonicalize observers and apply ConvertSmall; reapply it after every
operation, as described in Section~\ref{cse:small}. Use the ordinary BFR schedule: encapsulate and join when
both pipe ends are cutvertices; otherwise propagate a matched block-vertex
with a nontrivial embedding tree; otherwise simplify at the pipe mate of a
maximum-degree matched cutvertex, or at a matched vertex if none is a
cutvertex. The ordinary
applicability proof and potential depend on the graph and pipes, not on
the added frame and variable data~\cite{bfr}. The preceding updates retain
all required graph invariants. The subdivision vertices of PropagatePQ
are unmatched. Subdividing an edge changes neither whether an existing
vertex is a cutvertex nor the shapes of its embedding trees. A
subdivision vertex may itself be a cutvertex, when it subdivides a bridge;
this is harmless, because the BFR potential and the applicability of the
operations depend only on the pipes and on the degrees and cutvertex
status of their ends.

Let $m_0=|E(G)|$ and let $D_0$ be the largest initial pipe-end degree,
using zero when there are no pipes. There are at most $2m_0$ pipe
operations and no later pipe has degree greater than $D_0$. Each operation
adds $O(D_0)$ graph and marking data and at most $O(D_0^2)$ explicit
reference, category-domain, and relation-table data. Thus the size before
the terminal expansion is
\[
 S=O(M+m_0D_0^2)=O((M+1)^3).
\]
All embedding trees, marked subtrees, source constraints, and contour maps
can be computed by direct graph searches and SPQR operations. Even
traversing a whole current tree separately for every observer and listing
every injection gives a terminal system of size $O(S^2)$. The algebra
lemma and Gaussian elimination may conservatively be implemented in
cubic time in that size. Recomputing the graph structure at every
operation and executing every inverse map is bounded by a polynomial
of smaller degree. These estimates give $O((M+1)^{20})$, including
projection to the old Boolean variables and witness reconstruction.

The inverse records explicitly retain reference maps, frame equations,
detached-component categories, superblock intervals, and tree-contour
maps. Composing the exact preservation contracts proves
Theorem~\ref{cse:main}. Only explicit finite domains are enumerated in the
algebra lemma; no permutation family or planar-embedding family is
enumerated.

\paragraph{Finite adversarial checks.}
The accompanying exploratory script checked 35,240 instances of the
injective contour identity on four tree topologies with five or six
leaves, all their rotations, all proper anchor subsets of size at least
two, both admissible frames, and every remaining observed leaf. It also
compared the self-pipe category criterion with all cyclic orders for 720
anchor-preserving permutations on three through seven darts. These
checks passed. They support the specific identities tested and do not
replace the proofs or constitute an implementation of the full engine.
